% Replacement for the complete subsection "The surviving mixed relation"
% in chapters/04-shuffle-parity.tex. Do not insert in addition to that
% subsection: the legacy labels are intentionally retained.
% Add bibliography-entry.tex to the central bibliography.
\subsection{The mixed Gaussian relation: an exact Cayley certificate}

For the harmonic sums
\[
 S_p=\sum_{n\ge0}\frac{(-1)^nH_n}{(2n+1)^p},
\]
one has the exact bridge
\begin{equation}
 S_p=\Im\bigl(F_{p,1}(i)+\Li_{p,1}(i,-1)\bigr).
 \label{gaussian:eq:S-mixed}
\end{equation}
For an odd outer index $m=2n+1$, the inner coefficient
$\sum_{k<m}(1+(-1)^k)/k$ is $H_n$; even outer indices have zero
imaginary coefficient. This proves the bridge by absolute convergence
for $p>1$ and by convergent Abel limits for $p=1$.

\begin{theorem}[Mixed Gaussian reduction; exact Cayley certificate]
\label{gaussian:conj:S4}
\begin{equation}
 S_4=\frac{58g_{41}+24g_{32}}7
          +\frac{19\pi^5}{3584}-2\beta(4)\log2.
 \label{gaussian:eq:S4-short}
\end{equation}
\end{theorem}
\begin{proof}
The complete finite certificate and two independent exact verifiers are
supplied in \cite{CayleyS4Continuation2026}. Here is the analytic and
algebraic content of that certificate. Use the outer-letter-first forms
$Z=dt/t$ and $e_j=i^jdt/(1-i^jt)$, with conjugation fixing $Z,e_0,e_2$
and exchanging $e_1,e_3$. Put
\[
 A=Z^3e_1e_1,\quad B=Z^2e_1Ze_1,\quad C=Z^3e_1e_3,
 \quad V=Z^3e_1,\quad q=e_1-e_3.
\]
Write $\Delta P=P-\overline P$ and let $\mathbin{\sqcup\!\sqcup}$
denote shuffle multiplication. The target polynomial is
\[
 T=56\Delta C-408\Delta A-192\Delta B
       -19q^{\mathbin{\sqcup\!\sqcup}5}
       -112(\Delta V\mathbin{\sqcup\!\sqcup}e_2).
\]
Its iterated integral is
\[
 I(T)=16i\left(7S_4-58g_{41}-24g_{32}
                 -\frac{19\pi^5}{512}+14\beta(4)\log2\right).
\]
Indeed $I(q)=i\pi/2$, $I(e_2)=-\log2$, and
$I(\Delta V)=2i\beta(4)$.

The rational certificate expresses the odd-conjugation projection of
$T$ as 852 raw double-shuffle rows plus 224 times the odd projection of
$A-\mathcal D(A)$. Here $\mathcal D$ reverses word order and substitutes
\[
 Z\mapsto e_0-e_2,\quad e_0\mapsto Z+e_2,\quad
 e_1\mapsto e_2-e_3,\quad e_2\mapsto e_2,\quad
 e_3\mapsto e_2-e_1.
\]
The substitution $t\mapsto(1-t)/(1+t)$ proves $I(A)=I(\mathcal D(A))$;
all transformed words are admissible. Of the double-shuffle rows, 123
contain one initially divergent factor $\Li_1(1)$. Their divergent word
cancels, and equality of their cutoff and Abel finite parts is proved
in the cited continuation. Every other row uses only convergent values.
Thus every row has zero imaginary integral, and the verified rational
word identity gives $\Im I(T)=0$. Dividing by 16 and then by 7 proves
the formula. No numerical relation search is used as a proof premise.
\end{proof}

Equivalently, using the previously proved shuffle elimination,
\begin{equation}
 S_4=\frac{4g_{41}-3g_{32}-9g_{23}}7+\frac{\pi^5}{224}
        -\frac{27G\zeta(3)}{224}-2\beta(4)\log2.
 \label{gauss:eq:S4-closed}
\end{equation}
Both versions are now proved. The same certificate and one stuffle give
\begin{align}
 \Im\bigl(\Li_{4,1}(i,-1)-\Li_{1,4}(-1,i)\bigr)
 =\frac{102g_{41}+48g_{32}}7+\frac{79\pi^5}{10752}
     -3\beta(4)\log2.
 \label{cayley:eq:mixed-antisymmetric}
\end{align}
This is a weight-five mixed antisymmetric reduction to the $(i,1)$
double basket and single-value products. It does not prove minimal
depth or arithmetic independence of any remaining coordinate.
